\section{Target-based learning bounds at every realizable width}
\label{sec:critical-all-width-target-learning}

The quantitative target-mass constructions remove the width-limit
qualification from several learning bounds. Throughout this section,
training inputs are iid from any distribution $D$ on the Boolean cube,
labels are noiseless, and the learner is the original critical Gaussian
prior conditioned on a global absolute output margin and exact training
signs. A conclusion at \emph{every realizable width} means every integer
$N\ge1$ for which $\mu_N(\mathcal A_{N,f})>0$. Equivalently, a depth-$L$,
width-$N$ network of the specified architecture realizes the signs of $f$
with no zero output on the cube: rescaling its readout gives strict margin
$1/2$, and continuity makes that event have positive Gaussian probability.

Every sample bound below guarantees, with sample probability at least
$1-\delta$,
\[
 D\{x:p_N(f(x)\mid x)\le2/3\}\le\varepsilon,
 \qquad 0<\varepsilon,\delta<1.
\]
All constants are universal. The mild logarithmic accuracy factor comes
from the intrinsic finite-width theorem; no efficient posterior algorithm
is asserted.

\begin{lemma}[From a quantitative target mass to every-width learning]
\label{lem:critical-mass-to-all-width-learning}
Suppose $T\ge\log(d+L+2)$, $L\ge T$, and for a universal $c_0$ the
raw target cost obeys
\[
 a_N(f)\le V_*+c_0T^2/L\qquad(N\ge e^{c_0T}).
\]
Then, at every realizable width, a sufficient sample size is
\[
 n\ge\frac C\varepsilon
   \left\{e^{CT}(1+\log(1/\varepsilon))^c+\log(1/\delta)\right\}.
\]
\end{lemma}

\begin{proof}
Above the displayed width threshold, apply
Theorem~\ref{thm:critical-intrinsic-finite-width-learning}.
The target cost is at most $2+c_0T$, and its exponential factor is at
most $e^{CT}$. Since $d+L\le e^T$, all remaining resource factors are
absorbed in the asserted bound.

Below that threshold, the whole class of depth-$L$, width-$N$ tanh sign
networks has VC dimension polynomial in $d,L,N$, hence at most $e^{CT}$.
Sauer's lemma on the finite cube gives at most $\exp(e^{CT})$ distinct
sign functions after increasing $C$. A finite-class consistency bound
therefore makes every posterior-supported function have test error at
most $\varepsilon/3$ with the required sample probability. Averaging
and the probability-$2/3$ criterion give the result. Positive target
mass ensures that the conditioning event is nonempty for every sample.
\end{proof}

\begin{theorem}[Every-width learning from Fourier variation]
\label{thm:critical-all-width-fourier-learning}
Put $A_F(f)=\sum_S|\widehat f(S)|$ and $R=(d+L+2)A_F(f)$.
If $L\ge C\log R$, then at every realizable width a sufficient sample
size is
\[
 n\ge\frac C\varepsilon
 \left\{R^c(1+\log(1/\varepsilon))^c+\log(1/\delta)\right\}.
\]
In particular this is polynomial for every signed parity, independently
of its degree. The same conclusion holds for a decision tree with $s$
leaves or a length-$t$ binary-state branching program, with $R$ replaced
by $d+L+s+2$ or $d+L+t+2$, respectively. The latter programs may query
different coordinates in different states and may repeat queries.
\end{theorem}

\begin{proof}
Corollary~\ref{cor:critical-finite-fourier-mass} supplies the preceding
lemma with $T$ a universal multiple of $\log R$; enlarge the depth
constant accordingly. Every signed parity has $A_F=1$, while
Proposition~\ref{prop:critical-binary-state-learning} gives $A_F\le s$
for the tree and $A_F\le t+1$ for the binary-state program. Products
such as $(d+L+2)(t+1)$ are bounded by a universal power of the stated
replacement resource.
\end{proof}

\begin{theorem}[Every-width learning from a bounded threshold implementation]
\label{thm:critical-all-width-threshold-learning}
Suppose $f$ has a layered Boolean threshold implementation of depth
$h\ge1$, width at most $C_T$ including constant wires, unit gate margins,
and gate coefficient Euclidean norms at most $B\ge1$. Put
\[
 T=\log(d+2)+h\log(C_T+B+h+L+2).
\]
If $L\ge C T$, then at every realizable width a sufficient sample size is
\[
 n\ge\frac C\varepsilon
   \left\{e^{CT}(1+\log(1/\varepsilon))^c+\log(1/\delta)\right\}.
\]
In particular, the conclusion holds for every $N\ge C_T$.
\end{theorem}

\begin{proof}
The depth hypothesis implies the remaining-depth requirement
$L-h\ge C_0T$ in Corollary~\ref{cor:critical-finite-circuit-mass},
after increasing its universal constant. That corollary gives the
raw target-cost hypothesis of
Lemma~\ref{lem:critical-mass-to-all-width-learning}.

For the final assertion, multiply the coefficients of every circuit
gate by one sufficiently large common gain and replace its sign by
tanh. The unit margins and finitely many gates ensure that all hidden
activations uniformly approximate the corresponding Boolean wires
closely enough to preserve every later gate sign. Constant wires are
included in this approximation. Thus $h$ hidden layers of width $C_T$
and a linear readout realize the target. Any further hidden layers may
propagate its scalar output through a positive multiple of tanh,
preserving its sign. Rescale the final readout and pad unused neurons
with zero weights. This gives a strict-margin realization at every
$N\ge C_T$ and $L\ge h+1$, as ensured by the stated depth condition.
\end{proof}

\begin{corollary}[Real threshold teachers without a weight bound]
\label{cor:critical-all-width-real-threshold-learning}
Suppose $f$ has a layered real threshold circuit of depth $h\ge1$
and width $C_T\ge1$, with arbitrary coefficients and biases and a fixed
tie convention at each gate. Put $R=d+C_T+h+L+2$. If
$L\ge Ch\log R$, then at every realizable width a sufficient sample
size is
\[
 n\ge\frac C\varepsilon\left\{
 R^{ch}(1+\log(1/\varepsilon))^c+\log(1/\delta)\right\}.
\]
There is no dependence on the original coefficient magnitudes or
precisions. In particular, every Boolean halfspace has a universal
polynomial sample bound at every width $N\ge1$, whenever
$L\ge C\log(d+L+2)$.
\end{corollary}

\begin{proof}
Lemma~\ref{lem:critical-weight-independent-threshold-compiler} gives a
unit-margin implementation of depth $O(h)$, width and gate norms bounded
by a universal polynomial in $d+C_T+h+2$. Apply
Theorem~\ref{thm:critical-all-width-threshold-learning}; its parameter
is at most $Ch\log R$.
For a halfspace, perturb its bias if necessary to obtain strict signs
on the finite cube. A width-one tanh network realizes those signs at
any depth: use the affine form in the first layer and positive scalar
weights in each subsequent hidden layer, then scale the readout.
Thus every width is realizable.
\end{proof}

\begin{corollary}[Bounded-weight tanh teachers]
\label{cor:critical-all-width-analog-teacher-learning}
Suppose a depth-$D$, width-$C_T$ tanh teacher has parameter magnitudes
at most $B\ge1$ and output margin $1/2$ on the cube. Put
$R=d+C_T+B+D+L+2$. If $L\ge C(D+1)\log R$, then at every realizable
width a sufficient sample size is
\[
 n\ge\frac C\varepsilon
 \left\{R^{c(D+1)}(1+\log(1/\varepsilon))^c+\log(1/\delta)\right\}.
\]
\end{corollary}

\begin{proof}
Lemma~\ref{lem:analog-teacher-threshold-compilation} gives a threshold
circuit of depth $O(D+1)$, with width and gate coefficient norms bounded
by a universal polynomial in $d+C_T+B+D+2$. Its arithmetic circuits
use unweighted majority gates: the required precision does not introduce
a polynomial degree depending on $D$. The preceding theorem therefore
has $T\le C(D+1)\log R$ and proves the claim.
\end{proof}

\begin{corollary}[Low-rank teachers and solvable-group programs]
\label{cor:critical-all-width-structured-targets}
The following two classes have every-width bounds.
\begin{enumerate}
\item Suppose $f$ has a unit-margin tanh representation with product of
layer Frobenius norms at most $P\ge1$, and first-layer rank at most $r\ge1$.
Put $S=1+\log(d+L+2)+r\log(P+r+2)$. If $L\ge CS$, a sufficient sample
size at every realizable width is
\[
 \frac C\varepsilon\left\{
 (d+L+2)^c(P+r+2)^{cr}(1+\log(1/\varepsilon))^c
 +\log(1/\delta)\right\}.
\]
The approximate-rank condition in
Corollary~\ref{cor:critical-approx-low-rank-targets} has the same
consequence with $P$ replaced by $8P/7$.
\item Suppose $f$ is a length-$t$ program over a solvable group of
order $q\ge2$ and derived length $\ell\ge1$. Put $R=d+q+t+L+2$.
If $L\ge C\ell\log R$, a sufficient sample size at every realizable
width is
\[
 \frac C\varepsilon\left\{
 R^{c\ell}(1+\log(1/\varepsilon))^c+\log(1/\delta)\right\}.
\]
This includes a universal polynomial bound for permutation branching
programs of width at most four, independently of their sequential length.
\end{enumerate}
\end{corollary}

\begin{proof}
The proof of Theorem~\ref{thm:critical-low-rank-targets} constructs a
three-stage threshold circuit. Its width and gate norms are polynomial
in $P,r$ and the number of boxes $M\le(CP^2\sqrt r)^r$. The parameter
in Theorem~\ref{thm:critical-all-width-threshold-learning} is thus at
most $CS$. This proves the first assertion and its approximate-rank
variant. For the second,
Lemma~\ref{lem:critical-solvable-group-compiler} gives depth $O(\ell)$
and universal polynomial width and gate norms in $q,t$. The same
parameter is at most $C\ell\log R$. The permutation corollary uses
the bounded order and derived length of $S_4$, as before.
\end{proof}

\begin{theorem}[Every-width learning after Boolean layer blocking]
\label{thm:critical-all-width-boolean-blocking}
Suppose $f$ has a layered fan-in-two AND/OR/NOT circuit of width $C_T$
and depth $h\ge1$. For $1\le b\le h$, put $k=\lceil h/b\rceil$ and
\[
 U_b=\log(d+2)+k\{\log(C_T+h+L+2)+2^b\}.
\]
If $L\ge CU_b$, a sufficient sample size at every realizable width is
\[
 n\ge\frac C\varepsilon\left\{
 e^{CU_b}(1+\log(1/\varepsilon))^c+\log(1/\delta)\right\}.
\]
Consequently, for a fan-in-two circuit of size $s$, depth $h$, and
$S=\max\{16,d+s+2\}$, the conditions
\[
 L\ge C\log(S+L)
                \left(1+\frac h{\log\log(S+L)}\right),
\]
\[
 n\ge\frac C\varepsilon\left\{
 (S+L)^{C(1+h/\log\log(S+L))}
 (1+\log(1/\varepsilon))^c+\log(1/\delta)\right\}
\]
suffice at every realizable width. For formulas, $h$ may be replaced
by $C\log S$ after changing universal constants. At sufficiently large
learner depths polynomial in $S$, this is
$\exp\{O(\log^2S/\log\log S)\}$ apart from statistical factors.
\end{theorem}

\begin{proof}
The block construction in Theorem~\ref{thm:critical-boolean-blocking}
has $2k$ threshold stages, width at most
$(C_T+1)2^{C2^b}$, and coefficient norms at most $2^{C2^b}$ up to
universal factors. Its finite-width mass parameter is at most $CU_b$,
so apply Theorem~\ref{thm:critical-all-width-threshold-learning}.
For a size-$s$ circuit insert identity gates along skipped wires,
then use the same block choice and formula-balancing argument as in
Corollary~\ref{cor:critical-formula-blocking}. This yields
$U_b\le C\log(S+L)(1+h/\log\log(S+L))$.
\end{proof}

The parameter $U_b$ records the finite-width construction proved here.
We have not asserted that the sharper optimized-prefix parameter $T_b$
or $T_*$ in Section~\ref{sec:critical-optimized-blocking} has the same
finite-width target-mass estimate. The every-width formula bound does
retain the stated quasipolynomial order. None of these conclusions proves
joint polynomial sample complexity for arbitrary small deep circuits.
